\begin{afterword}
\summaryhead

The post opens with the scene in Orwell's \textsc{Nineteen Eighty-Four} in which O'Brien burns a photograph and says he does not remember it, and Winston wonders whether such a ``lunatic dislocation in the mind'' could really happen. The post then asks whether it would be wise to choose to stay biased if bias made us happier, and the author doubts that it could work.

We cannot make ourselves believe things by an act of will. A rational choice to be biased would need an accurate picture of the consequences, which would then have to be forgotten, and the forgetting forgotten. Without that picture, staying biased is a blind choice, ``willful stupidity.'' An optimistic driver may skip the seatbelt happily for years and then crash, and people neglect such risks because unknown dangers do not come to mind.

The author grants that stupid people may be happier, but holds that readers who know about biases cannot return to ignorance. They should aim for the happiness a rationalist can achieve, which may prove greater.

\responsehead

The post argues in two steps that deliberately choosing a bias cannot work, and each step covers fewer cases than the conclusion.

The first step is that we cannot believe at will: ``You cannot make yourself believe the sky is green by an act of will.'' This answers only direct control. The old proposal for acquiring a belief on purpose is indirect. Pascal told the unbeliever who says ``I am so made that I cannot believe'' to do what believers had done at first, acting ``as if they believed,'' and said that this ``will naturally make you believe.'' The woman in ``No, Really, I've Deceived Myself'' says she took that route. The post does not discuss it. The author also later described the act-of-will sentence as partly ``trying to instill a self-fulfilling prophecy,'' not only a report of fact.

The second step is that a rational choice to be biased needs the very knowledge the bias would remove. For a single belief chosen with full knowledge of its effects, this is a good argument, and it is Orwell's own account of doublethink, which the post credits. But the post turns it into a rule that ``You can't know the consequences of being biased, until you have already debiased yourself,'' and allows only two cases: full knowledge, or a choice made ``without any clear idea of the consequences.'' Partial knowledge from studies of other people, the kind of knowledge every decision uses, is left out.

That omission matters because the question the post opens with has a research literature. Taylor and Brown (1988) argued that positive illusions, among them the belief that one drives better than others, are normal and appear to promote mental health and happiness. Their thesis was contested, but it is evidence for the view the post rejects. The post names no one who holds that view and answers it with a driving case in which the bias leads to a dangerous act, not with the cases of threat and bad news where Taylor and Brown suggested illusions may help most.

The ending moves from ``I suspect the vast majority'' of readers could not achieve ignorance to ``You can never achieve that degree of ignorance,'' with no evidence for either.

In short: an argument against choosing to be biased that answers only belief by a direct act of will and choices made with full knowledge or none, and does not engage the research that raises the question.
\end{afterword}
